\documentclass[11pt,a4paper]{article}

\usepackage{amsmath,amssymb,amsthm}
\usepackage{mathtools}
\usepackage{hyperref}
\usepackage{cleveref}
\usepackage[margin=1in]{geometry}
\usepackage{enumitem}
\usepackage{booktabs}
\usepackage{bussproofs}

\newtheorem{theorem}{Theorem}[section]
\newtheorem{lemma}[theorem]{Lemma}
\newtheorem{proposition}[theorem]{Proposition}
\newtheorem{corollary}[theorem]{Corollary}
\newtheorem{definition}[theorem]{Definition}
\newtheorem{axiom}{Axiom}
\newtheorem{remark}[theorem]{Remark}
\newtheorem{example}[theorem]{Example}

\newcommand{\Prob}{\mathsf{Prob}}
\newcommand{\ProbProp}{\mathsf{ProbProp}}
\newcommand{\E}{\mathbb{E}}
\newcommand{\pand}{\wedge_p}
\newcommand{\por}{\vee_p}
\newcommand{\pnot}{\neg_p}
\newcommand{\pzero}{\mathbb{0}}
\newcommand{\pone}{\mathbb{1}}
\newcommand{\ple}{\leq_p}
\newcommand{\padd}{+_p}
\newcommand{\pmul}{\cdot_p}
\newcommand{\psub}{-_p}
\newcommand{\entails}[1]{\vdash_{#1}}

\title{Probabilistic Proof Theory\\[0.5em]
\large Sequent Calculus with Confidence Grades}

\author{Probabilistic Foundations Project}

\date{\today}

\begin{document}

\maketitle

\begin{abstract}
We develop a proof theory for probabilistic reasoning where derivations carry confidence grades: $\Gamma \entails{p} \phi$ means ``from $\Gamma$, derive $\phi$ with confidence at least $p$.'' The classical sequent calculus is recovered when $p = \pone$, but the interesting mathematics happens for $p < \pone$. We present four fundamental proof rules---axiom, weakening, cut, and monotonicity---and show how confidence degrades through inference chains. This framework supports reasoning under uncertainty as a first-class citizen, not an afterthought.
\end{abstract}

\section{Introduction}

\subsection{Classical vs. Probabilistic Entailment}

In classical logic, entailment is binary: either $\Gamma \vdash \phi$ or not. This forces a sharp boundary between ``known'' and ``unknown,'' with no room for partial confidence.

In probabilistic proof theory, entailment is graded:
\[
\Gamma \entails{p} \phi
\]
This means: ``The confidence in $\phi$ given $\Gamma$ is at least $p$.''

\subsection{Where the Interesting Mathematics Lives}

The classical case $p = \pone$ recovers standard sequent calculus. But the value of this framework is in the $p < \pone$ regime:

\begin{itemize}
    \item Confidence propagates through proofs: if $\Gamma \entails{p} \phi$ and $\phi \entails{q} \psi$, then $\Gamma \entails{p \pmul q} \psi$
    \item Approximate reasoning becomes formal: ``likely true'' has precise meaning
    \item Bounds are compositional: complex derivations have calculable confidence
\end{itemize}

\subsection{Prerequisites}

This paper assumes familiarity with Paper 1 (the $\Prob$ type and $\ProbProp$). We introduce conditional expectation $\E[\cdot|\cdot]$ as a new primitive.

\section{Conditional Expectation}

\subsection{The Core Primitive}

Conditional expectation is the fundamental operation for reasoning about one proposition given another.

\begin{axiom}[Conditional Expectation]
There exists an operation:
\[
\E[\cdot|\cdot] : \ProbProp(\alpha) \times \ProbProp(\alpha) \to \Prob
\]
written $\E[\phi|\Gamma]$, representing the expected value of $\phi$ given $\Gamma$.
\end{axiom}

\subsection{Axioms for Conditional Expectation}

\begin{axiom}[Normalization]
$\E[\mathbf{1}|\Gamma] = \pone$ where $\mathbf{1}(x) = \pone$ for all $x$.
\end{axiom}

\begin{axiom}[Monotonicity]
If $\phi(x) \ple \psi(x)$ for all $x$, then $\E[\phi|\Gamma] \ple \E[\psi|\Gamma]$.
\end{axiom}

\begin{axiom}[Self-Conditioning]
$\E[\phi|\phi] = \pone$.
\end{axiom}

This captures: if we condition on $\phi$, the expected value of $\phi$ is maximal.

\begin{axiom}[Linearity]
$\E[c \pmul \phi|\Gamma] = c \pmul \E[\phi|\Gamma]$ for constant $c : \Prob$.
\end{axiom}

\begin{axiom}[Additivity]
$\E[\phi \padd \psi|\Gamma] = \E[\phi|\Gamma] \padd \E[\psi|\Gamma]$.
\end{axiom}

\begin{axiom}[Product Rule]
$\E[\phi \pmul \psi|\chi] = \E[\phi|\psi \pmul \chi] \pmul \E[\psi|\chi]$.
\end{axiom}

This is Bayes' theorem in expectation form.

\begin{axiom}[Chain Rule / Tower Law]
$\E[\psi|\chi] \pmul \E[\phi|\psi] \ple \E[\phi|\chi]$.
\end{axiom}

Conditioning through an intermediate proposition can only decrease confidence.

\subsection{Interpretation}

Think of $\E[\phi|\Gamma]$ as: ``How much do I expect $\phi$ to be true, given that $\Gamma$ describes my evidence?''

For classical (0/1-valued) propositions, this reduces to conditional probability: $\E[\phi|\Gamma] = P(\phi|\Gamma)$.

For general probabilistic propositions, it is the expected value of the function $\phi$ under the conditional distribution.

\section{Probabilistic Entailment}

\subsection{Definition}

\begin{definition}[Probabilistic Entailment]
For $\Gamma, \phi : \ProbProp(\alpha)$ and $p : \Prob$, we define:
\[
\Gamma \entails{p} \phi \quad :\Leftrightarrow \quad p \ple \E[\phi|\Gamma]
\]
\end{definition}

This means: ``The expected value of $\phi$ given $\Gamma$ is at least $p$.''

\begin{remark}
Higher $p$ is a stronger claim. $\Gamma \entails{\pone} \phi$ is the strongest: perfect confidence.
\end{remark}

\subsection{Classical Entailment}

\begin{definition}[Classical Entailment]
For classical propositions $\Gamma, \phi$:
\[
\Gamma \vdash^c \phi \quad :\Leftrightarrow \quad \Gamma \entails{\pone} \phi
\]
\end{definition}

The classical sequent $\Gamma \vdash \phi$ is the special case of probabilistic entailment with maximal confidence.

\section{Proof Rules}

\subsection{The Four Fundamental Rules}

\begin{theorem}[Axiom Rule]
For any $\phi : \ProbProp(\alpha)$:
\[
\phi \entails{\pone} \phi
\]
\end{theorem}

\begin{proof}
By Self-Conditioning: $\E[\phi|\phi] = \pone$, so $\pone \ple \E[\phi|\phi]$.
\end{proof}

\begin{theorem}[Weakening]
If $\Gamma \entails{p} \phi$ and $q \ple p$, then $\Gamma \entails{q} \phi$.
\end{theorem}

\begin{proof}
We have $p \ple \E[\phi|\Gamma]$. By transitivity of $\ple$: $q \ple p \ple \E[\phi|\Gamma]$.
\end{proof}

\begin{remark}
Weakening \emph{decreases} confidence. This is not ``weakening the hypothesis'' as in classical logic---it is ``accepting lower confidence.''
\end{remark}

\begin{theorem}[Cut Rule]
If $\Gamma \entails{p} \phi$ and $\phi \entails{q} \psi$, then $\Gamma \entails{p \pmul q} \psi$.
\end{theorem}

\begin{proof}
We have:
\begin{align*}
    p &\ple \E[\phi|\Gamma] \\
    q &\ple \E[\psi|\phi]
\end{align*}
By monotonicity of multiplication:
\[
p \pmul q \ple \E[\phi|\Gamma] \pmul \E[\psi|\phi]
\]
By the Chain Rule:
\[
\E[\phi|\Gamma] \pmul \E[\psi|\phi] \ple \E[\psi|\Gamma]
\]
By transitivity: $p \pmul q \ple \E[\psi|\Gamma]$.
\end{proof}

\begin{remark}
This is the key rule for confidence propagation. Each inference step \emph{multiplies} confidence. A chain of 10 steps, each with confidence 0.9, yields overall confidence $0.9^{10} \approx 0.35$.
\end{remark}

\begin{theorem}[Monotonicity Rule]
If $\phi(x) \ple \psi(x)$ for all $x$, and $\Gamma \entails{p} \phi$, then $\Gamma \entails{p} \psi$.
\end{theorem}

\begin{proof}
By monotonicity of conditional expectation: $\E[\phi|\Gamma] \ple \E[\psi|\Gamma]$.
Thus $p \ple \E[\phi|\Gamma] \ple \E[\psi|\Gamma]$.
\end{proof}

\subsection{Derived Rules}

\begin{corollary}[Transitivity]
If $\Gamma \entails{p} \phi$, $\phi \entails{q} \psi$, and $\psi \entails{r} \chi$, then:
\[
\Gamma \entails{p \pmul q \pmul r} \chi
\]
\end{corollary}

\begin{corollary}[Bottom Bound]
$\Gamma \entails{\pzero} \phi$ for any $\Gamma, \phi$.
\end{corollary}

\begin{proof}
$\pzero \ple p$ for all $p$, including $\E[\phi|\Gamma]$.
\end{proof}

\section{Structural Rules}

\subsection{Exchange}

\begin{theorem}[Exchange]
If $(A \pand B) \pand \Gamma \entails{p} C$, then $(B \pand A) \pand \Gamma \entails{p} C$.
\end{theorem}

\begin{proof}
Since $(A \pand B)(x) = A(x) \pmul B(x) = B(x) \pmul A(x) = (B \pand A)(x)$, the contexts are equal pointwise, so the conditional expectations are equal.
\end{proof}

\subsection{Contraction (Classical Only)}

\begin{theorem}[Contraction]
For classical $A$: if $(A \pand A) \pand \Gamma \entails{p} B$, then $A \pand \Gamma \entails{p} B$.
\end{theorem}

\begin{proof}
For classical $A$: $(A \pand A)(x) = A(x) \pmul A(x) = A(x)$ (since $\pzero \pmul \pzero = \pzero$ and $\pone \pmul \pone = \pone$).
\end{proof}

\begin{remark}
Contraction \emph{fails} for non-classical propositions. If $A(x) = 0.5$, then $(A \pand A)(x) = 0.25 \neq 0.5$. This is a feature, not a bug: in probabilistic reasoning, using evidence twice is not the same as using it once.
\end{remark}

\subsection{Associativity}

\begin{theorem}[Associativity]
$((A \pand B) \pand C) \pand \Gamma \entails{p} D$ iff $(A \pand (B \pand C)) \pand \Gamma \entails{p} D$.
\end{theorem}

\begin{proof}
Multiplication in $\Prob$ is associative.
\end{proof}

\section{Example Derivations}

\subsection{Confidence Degradation Chain}

Suppose:
\begin{itemize}
    \item $A \entails{0.9} B$ (``$A$ strongly suggests $B$'')
    \item $B \entails{0.8} C$ (``$B$ moderately suggests $C$'')
    \item $C \entails{0.7} D$ (``$C$ weakly suggests $D$'')
\end{itemize}

By repeated application of Cut:
\[
A \entails{0.9 \times 0.8 \times 0.7} D = A \entails{0.504} D
\]

Three ``reasonable'' inferences combine to give only moderate overall confidence.

\subsection{Comparing to Classical Logic}

In classical logic, if $A \vdash B$, $B \vdash C$, $C \vdash D$, then $A \vdash D$ with no degradation.

This corresponds to the special case where all confidences are $\pone$:
\[
A \entails{\pone} B, \quad B \entails{\pone} C, \quad C \entails{\pone} D \quad \Rightarrow \quad A \entails{\pone} D
\]

Classical logic is the idealized limit where every step is perfectly reliable.

\section{The LEM Expectation Identity}

\begin{theorem}[LEM for Expectations]
For any $\phi, \psi : \ProbProp(\alpha)$:
\[
\E[\phi|\psi] \padd \E[\pnot\phi|\psi] = \pone
\]
\end{theorem}

\begin{proof}
By additivity:
\[
\E[\phi \padd (\pone \psub \phi)|\psi] = \E[\phi|\psi] \padd \E[\pone \psub \phi|\psi]
\]
Since $\phi(x) \padd (\pone \psub \phi(x)) = \pone$ for all $x$, the LHS equals $\E[\mathbf{1}|\psi] = \pone$.
\end{proof}

\begin{remark}
This is not ``LEM holds in expectation.'' It is: expectations of a proposition and its negation always sum to 1. The logical law is a consequence of probability algebra.
\end{remark}

\section{Negation}

\subsection{Weak Negation}

Weak negation is pointwise complement:
\[
(\pnot\phi)(x) = \pone \psub \phi(x)
\]

\subsection{Strong Negation}

\begin{definition}[Strong Negation]
A proposition $\phi$ is \emph{strongly negated} if $\E[\phi|\psi] = \pzero$ for all $\psi$.
\end{definition}

This means: $\phi$ has probability zero under every conditioning.

\begin{theorem}
If $\phi$ is strongly negated, then $\E[\pnot\phi|\psi] = \pone$ for all $\psi$.
\end{theorem}

\begin{proof}
By LEM: $\E[\phi|\psi] \padd \E[\pnot\phi|\psi] = \pone$. Since $\E[\phi|\psi] = \pzero$, we get $\E[\pnot\phi|\psi] = \pone$.
\end{proof}

\subsection{Strong Negation Propagates}

\begin{theorem}
If $\phi$ is strongly negated, so is $\phi \pand \psi$ for any $\psi$.
\end{theorem}

\begin{proof}
$(\phi \pand \psi)(x) = \phi(x) \pmul \psi(x) \ple \phi(x) \pmul \pone = \phi(x)$.
By monotonicity: $\E[\phi \pand \psi|\chi] \ple \E[\phi|\chi] = \pzero$.
\end{proof}

\section{Implication}

\subsection{Definition}

\begin{definition}[Probabilistic Implication]
$(A \to_p B)(x) := \pone \psub (A(x) \pmul (\pone \psub B(x)))$
\end{definition}

\subsection{Properties (Classical Case)}

\begin{theorem}[Reflexivity]
For classical $A$: $(A \to_p A)(x) = \pone$.
\end{theorem}

\begin{proof}
If $A(x) = \pzero$: $\pone \psub (\pzero \pmul (\pone \psub \pzero)) = \pone \psub \pzero = \pone$.
If $A(x) = \pone$: $\pone \psub (\pone \pmul (\pone \psub \pone)) = \pone \psub \pzero = \pone$.
\end{proof}

\begin{theorem}[Modus Ponens Bound]
For classical $A, B$: $A(x) \pmul (A \to_p B)(x) \ple B(x)$.
\end{theorem}

\begin{proof}
Case analysis on $A(x), B(x) \in \{\pzero, \pone\}$.
\end{proof}

\begin{remark}
These are exact equalities only for classical propositions. For general probabilistic propositions, implication has more nuanced behavior.
\end{remark}

\section{Implementation}

\subsection{Lean 4 Code}

\begin{verbatim}
axiom CondExp {a : Type} : ProbProp a -> ProbProp a -> Prob

notation "E[" f " | " g "]" => CondExp f g

structure ProbEntails {a : Type} (G f : ProbProp a) (p : Prob) where
  bound : p <= E[f | G]

notation:25 G " |-[" p "] " f => ProbEntails G f p

theorem axiom_rule {a : Type} (f : ProbProp a) : f |-[1] f := by
  constructor
  rw [CondExp.self_norm]
  exact Prob.le_refl 1

theorem cut_rule {a : Type} (G f y : ProbProp a) (p q : Prob) :
    (G |-[p] f) -> (f |-[q] y) -> (G |-[p * q] y) := by
  intro h1 h2
  constructor
  have hp : p <= E[f | G] := h1.bound
  have hq : q <= E[y | f] := h2.bound
  have hmul : (p * q) <= (E[f | G] * E[y | f]) :=
    Prob.mul_le_mul p q _ _ hp hq
  have hchain : (E[f | G] * E[y | f]) <= E[y | G] :=
    CondExp.chain y f G
  exact Prob.le_trans (p * q) _ _ hmul hchain
\end{verbatim}

\subsection{Verification Status}

All proof rules are machine-verified in Lean 4.

\section{Discussion}

\subsection{Why Not Just Use Probability Theory?}

Standard probability theory is built on measure theory, which is built on set theory, which is built on classical logic. Our approach inverts this: conditional expectation is primitive, and classical logic is a derived special case.

The advantage is not that we get ``more probability theory''---it is that we get a different foundation where uncertainty is native.

\subsection{Relation to Fuzzy Logic}

Fuzzy logic also has graded truth values, but uses different operations (typically $\min/\max$ for $\land/\lor$). Our operations come from probability theory (product for conjunction, inclusion-exclusion for disjunction).

\subsection{What Comes Next}

\begin{itemize}
    \item \textbf{Paper 3}: Natural numbers as distributions
    \item \textbf{Paper 4}: Existence proofs via distributions (no AC)
    \item \textbf{Paper 5}: Synthetic construction of $\Prob$
\end{itemize}

\section{Conclusion}

Probabilistic proof theory extends classical sequent calculus with confidence grades. The key insight is that confidence propagates multiplicatively: each inference step can only decrease overall confidence.

The classical case $p = \pone$ is the idealized limit of perfect reliability. The interesting mathematics happens when $p < \pone$, where we can formally reason about uncertainty.

\bibliographystyle{plain}

\end{document}
